\scanpage{127}
\mspage{42}
% archive sheet number: 126
\margin{\deleted{$(\unsure{m=q(q+1)|m-q(q+1)})$}}

% L01-04
(ii) Il existe un sous-groupe \inserted{ouvert} $I'$ de $I$, \deleted{et un entier \unsure{$m$}} tel que pour tout $\nu$, l'action de $I'$ sur $(\underline{\Psi}_\nu^q)_{\mathrm{tame}}$ soit unipotente d'échelon \deleted{$\unsure{q(q+1)/2}$} $q+1$. \deleted{Alors}\inserted{Alors} L'action de $I'$ sur $(\underline{\Psi}^q)_{\mathrm{tame}}$ est également unipotente d'échelon \deleted{$\unsure{q(q+1)/2}$} $q+1$.

% L05-06
\uline{Démonstration} Soient $g,g'$ les morphismes canoniques de $U_{\overline\eta_t}$ dans $X',X$. On a

% L06 left margin
\diagram{\begin{tikzcd}[ampersand replacement=\&]
X'\arrow[d,"p"']\&\\
X\&U_{\overline\eta_t}\deleted{\unsure{(K_t)}}\arrow[l,"g"']\arrow[ul,"g'"']
\end{tikzcd}}

% L07
\[
(\underline{\Psi}_\nu^q)_{\mathrm{tame}}=R^q g_*^{\deleted{\unsure{\prime|2}}}(\Lambda_\nu)|X_0
\tag*{(6.2.1)}
\]

% L08-09
et \deleted{par suite} d'autre part la suite spectrale de Leray du morphisme composé $g=pg'$ donne

% L10
\[
R^*g_*(\Lambda_\nu)=(\underline{\Psi}_\nu^*)_{\mathrm{tame}}\Longleftarrow\deleted{\unsure{E_2^{pq}}}\inserted{E_2^{pq}(\nu)}=R^p p_*(R^q g'_*(\Lambda_\nu^{\deleted{\unsure{q}}})).
\tag*{(6.2.2)}
\]

% L11-13
En vertu des \deleted{\unsure{résultats}}\inserted{calculs du § 4}, on trouve un s-groupe ouvert $I'$ de $I$ qui, \inserted{pour tt $q,\nu$,} opère \deleted{de façon unipotente d'échelon \unsure{$q+1$} sur}\inserted{trivialement sur} $R^q g'_*(\underline{\Psi}_\nu^q)$. Il s'ensuit que $I'$ opère \deleted{\unsure{de}}\inserted{trivialement} \deleted{de façon unipotente d'échelon $\unsure{(0+1)+(1+1)+\cdots+(q+1)}$} sur $E_2^{pq}(\nu)$, donc \deleted{sur} il opère de façon unipotente \inserted{d'échelon $q+1$} \deleted{\unsure{dans}}

% L14-15
\deleted{\unsure{échelon} $(\underline{\Psi}_\nu^q)_{\mathrm{tame}}$, avec pour échelon \unsure{$q+1$}}
\[
\deleted{(0+1)+(1+1)+\cdots+(q+1)=\frac{q(q+1)}{2}},
\]

% L16
ce qui prouve la première assertion de ii).

% L17-20
D'autre part, les calculs explicites du § 4 montrent que les $R^q g'_*(\unsure{\Lambda|\underline{\Psi}}_\nu^{\deleted{q}})|X'_0$ sont des faisceaux \uline{constructibles}, qui pour $\nu$ variable forment un faisceau \deleted{\unsure{const.}} $\ell$-adique (i.e. $R^q g'_*(\Lambda_{\nu'})|X'_0=(R^q g'_*(\Lambda_\nu)|X'_0)\otimes_{\Lambda_\nu}\unsure{\Lambda_{\nu\nu'}|\Lambda_{\nu'}}$ si $\nu\geq\nu'$).

% L21
En vertu de SGA 5 \uline{VI} 2.2, il s'ensuit que
